\newpage
\pdfbookmark[0]{Kurzfassung}{kurzfassung}
\chapter*{Kurzfassung}
\markboth{Kurzfassung}{Kurzfassung}

Diese Dissertation befasst sich mit der erweiterten 2D- und 3D-Charakterisierung von Klinkerphasen und hydratisierten zementären Bindemitteln.
Herkömmliche Charakterisierungsmethoden stoßen bei der Auflösung von Nanoporosität oder der chemischen Sensitivität, um Spurenelemente bestimmten Klinkerphasen zuzuordnen, oft an ihre Grenzen.
Das Hauptziel dieser Arbeit ist es, diese Einschränkungen durch die Etablierung korrelativer Mikroskopie-Workflows und automatisierter Auswertungswerkzeuge zu überwinden, die die Lücke zwischen hochauflösender Bildgebung und der Bestimmung physikalischer Eigenschaften schließen.
Dazu werden verschiedene moderne bildgebende und analytische Verfahren wie die Rasterelektronenmikroskopie (REM), die fokussierte Ionenstrahl-Nano\-to\-mo\-gra\-phie (\acs*{FIBnt}), die Laserablation mit induktiv gekoppeltem Plasma und Massenspektrometrie (\acs*{LAICPMS}), die Argon-Ionenstrahlpolitur (\acs*{ArBIB}), die Tieftemperatur-Differenzkalorimetrie (\acs*{LTDSC}) und die Nanoindentation (\acs*{NI}) systematisch kombiniert.

Eine zentrale Säule dieser Forschung ist der Einsatz moderner Präparationstechniken.
\acs*{ArBIB} wird verwendet, um artefaktfreie Oberflächen zu erzeugen, was für die hochauflösende REM-Bildgebung entscheidend ist, da sie die native Porenstruktur ohne Harzeinbettung erhält.
Diese Methode erwies sich als unerlässlich für die Präparation von Proben für nachfolgende korrelative Analysen.

Der korrelative Ansatz wird anhand mehrerer Workflows demonstriert.
Die Kombination von Energiedispersiver Röntgenspektroskopie im REM mit \acs*{NI} oder \acs*{LAICPMS} ermöglicht die direkte Zuordnung von phasenspezifischen mechanischen Eigenschaften und Spurenelementverteilungen.
Dies geht über traditionelle Dekonvolutionsmethoden hinaus, indem chemische und mechanische Daten mit spezifischen mikrostrukturellen Merkmalen verknüpft werden, wie z.\,B. die Zuordnung der lokalen Konzentration von Spurenelementen zu einzelnen Klinkerphasen oder die Bestimmung der Härte einzelner Hydratphasen.

Um von der qualitativen Beobachtung von mikrostrukturellen Features zur quantitativen Analyse überzugehen, wurden automatisierte Bildanalysewerkzeuge entwickelt.
Ein neuartiger Algorithmus ermöglicht die automatisierte Messung der Hydratsaumdicke (\acs*{HLT}) und von Partikelabständen aus großflächigen 2D-Bildern.
Dies erleichtert die statistische Analyse der Hydratationskinetik, wie an \acs*{C3S}- und \acs*{C2S}-Systemen gezeigt, und quantifiziert die Dispergierwirkung von Zusatzmitteln oder physikalischen Behandlungen wie dem Hochleistungsultraschall.
Darüber hinaus deutete die Korrelation von \acs*{HLT}-Daten mit Messungen der Elektronenrückstreubeugung darauf hin, dass das Hydratsaumwachstum stärker von der lokalen räumlichen Verfügbarkeit und Porenkonnektivität als von der kristallographischen Orientierung des darunterliegenden Klinkerkorns bestimmt wird.

Für die 3D-Charakterisierung wurden optimierte Workflows für die \acs*{FIBnt} etabliert.
Eine verbesserte Lift-out-Technik für zementäre Materialien wurde implementiert, um Bildartefakte zu minimieren, und die Femtosekunden-Laserablation  wurde genutzt, um schnell größere, repräsentativere Volumina zu erschließen.
Diese Präparationsmethoden werden durch eine automatisierte Nachbearbeitungspipeline für 3D-Daten ergänzt, die Registrierung, Entrauschen und eine auf Superpixeln basierende Segmentierung zur Erzeugung quantitativer Phasenkarten umfasst.
Diese 3D-Datensätze wurden zusammen mit 2D-Bildgebung, \acs*{LTDSC} und der Quecksilberdruckporosimetrie in einer vergleichenden Studie von Gel- und Kapillarporosität verwendet, die eine lineare Korrelation zwischen der Reduzierung der Kapillarporosität und der Bildung von Gelporosität während der Hydratation aufzeigte.

Die Ergebnisse zeigen, dass ein korrelativer, multimodaler Ansatz ein umfassenderes Verständnis liefert, als es einzelne Techniken allein könnten.
Die entwickelten automatisierten Werkzeuge liefern objektive und statistisch robuste Daten und legen damit den Grundstein für zukünftige Forschung, die darauf abzielt, die Effizienz und Dauerhaftigkeit von Bindemitteln durch ein tieferes, datengestütztes Verständnis der zementären Mikrostruktur zu optimieren.

